\documentclass[dvipdfmx,11pt]{beamer}

\usetheme{Madrid}
\usecolortheme{default}
\setbeamertemplate{navigation symbols}{}

\usepackage{amsmath,amssymb,bm}
\usepackage{booktabs}

\newcommand{\T}{\mathcal{T}}
\newcommand{\R}{\mathbb{R}}
\newcommand{\1}{\mathbf{1}}

\title{重複留保カテゴリと点数制約の下での公平な admissions}
\subtitle{SPSE と PA-ODR}
\author{}
\date{}

\begin{document}

\begin{frame}
  \titlepage
\end{frame}

\begin{frame}{問題設定}
  一つの学校における admissions を考える。
  \begin{itemize}
    \item 一般カテゴリ $o$ と複数の留保カテゴリ $T$
    \item 学生 $i$ は複数の留保カテゴリに同時に資格を持ちうる
    （$|\T(i)|\ge 2$ のとき \textbf{overlapping}）
    \item \textbf{One-to-one counting}:
    合格者は高々一つのカテゴリにしか計上されない
    \\ $\Rightarrow$ 「誰を取るか」だけでなく「どのカテゴリで数えるか」自体が配分の一部
    \item \textbf{点数制約}: 一般 cutoff と各留保 cutoff の差を $\delta$ 以内に抑える
  \end{itemize}


\end{frame}

\section{Motivation}

\begin{frame}{モチベーション 1: 加点方式を重複留保へ拡張する}
  出発点は，留保カテゴリの学生に加点し，
  加点後スコアで同じ選抜に乗せる発想である。

  \vspace{0.4em}
  \begin{block}{例}
    女性に 10 点を加点して，男性と同じプールで競わせる。
  \end{block}

  \vspace{0.4em}
  これを重複カテゴリに拡張すると，
  \begin{itemize}
    \item 学生は複数の留保カテゴリに資格を持ちうる
    \item しかし one-to-one counting では，高々一つのカテゴリにしか計上できない
    \item したがって，重複資格学生を\textbf{どのカテゴリに置くか}が新しい選択になる
  \end{itemize}
\end{frame}

\begin{frame}{モチベーション 2: cutoff 差を制度制約として扱う}
  本稿では，score-feasibility を通常の raw score cutoff ではなく，
  \textbf{留保カテゴリ側では加点後 cutoff} を用いて定義する。

  \vspace{0.4em}
  \[
    \text{open cutoff は raw score},\qquad
    \text{reserved cutoff は elevated score}
  \]

  \vspace{0.4em}
  \begin{itemize}
    \item 加点は制度上正当な補正として扱う
    \item そのうえで，一般 cutoff と留保 cutoff の差を $\delta$ 以内に抑える
    \item 加点後スコアで DA を回すだけでは，この cutoff 差は保証されない
  \end{itemize}

  \vspace{0.4em}
  \begin{block}{本研究の問い}
    重複カテゴリと one-to-one counting の下で，
    score-feasibility を守るメカニズムをどう設計するか。
  \end{block}
\end{frame}

% \begin{frame}{本研究の貢献}
%   \begin{enumerate}
%     \item 共通加点と cutoff 差制約を持つ重複留保モデルを定式化し，
%     non-wastefulness は点数制約に相対化して述べ直す必要がある一方，
%     within-category fairness は述べ直す必要が\textbf{ない}ことを示す
%     \item SPSE（student-proposing score-elevated mechanism）は
%     feasible・non-wasteful・within-category fairness だが
%     score-feasible とは限らない。
%     その破れには 2 つの原因があり，\textbf{overlap 由来の方が新しい}
%     \item score-feasibility・open-completeness・reserve-completeness は，
%     \textbf{各留保カテゴリが個別には score-feasible に充足可能でも}同時には満たせない
%     （Hall 条件の破れ；overlap がなければ消える）
%     \item PA-ODR を提案: score-feasibility を hard constraint として
%     充足不能な留保容量を de-reserve。
%     within-category fairness は\textbf{弱めずに}保つが，
%     最小性を失う
%     \item 留保側の規則は 1 回の最大重み二部マッチングとして実現でき，
%     最小 de-reservation は $|T|$ が定数なら多項式時間で求まる
%   \end{enumerate}
% \end{frame}

\begin{frame}{モデル 1: 基本記号}
  学生集合 $I$，留保カテゴリ集合 $T$，一般カテゴリ $o$，
  \[
    V := T\cup\{o\},\qquad q^c\in\mathbb{N}\ \ (c\in V).
  \]

  各学生 $i\in I$ は
  \begin{itemize}
    \item 加点前スコア $\sigma(i)\in\R_+$（学生間で相異なる）
    \item 資格カテゴリ集合 $\T(i)\subseteq T$
    \item 応募可能集合 $V_i:=\{o\}\cup\T(i)$ 上の厳密選好 $\succ_i$
    \\ $r_i(c)$: $i$ にとって $c$ が何番目の希望か（$1$ が最良）
  \end{itemize}

  共通パラメータ:
  \[
    k\ge 0\ (\text{全留保カテゴリ共通の加点}),\qquad
    \delta\ge 0\ (\text{許容 cutoff 差}).
  \]

  % \vspace{0.3em}
  % {\small 学校は 1 つなので，学生は学校ではなく\textbf{入学ルート}を順序づける。}
\end{frame}

\begin{frame}{モデル 2: One-to-one counting と feasibility}
  合格者は資格を持つ全カテゴリではなく，
  メカニズムが選ぶ\textbf{ちょうど一つ}のカテゴリに計上される。

  \begin{definition}[Matching と feasibility]
    $\mu=(\mu_c)_{c\in V}$，$M(\mu):=\bigcup_{c\in V}\mu_c$。
    $\mu$ が \textbf{feasible} とは
    \begin{enumerate}
      \item[(i)] $|\mu_c|\le q^c\quad(\forall c\in V)$
      \item[(ii)] $\mu_c\cap\mu_{c'}=\emptyset\quad(\forall c\neq c')$
      \item[(iii)] $c\in V_i\quad(\forall i\in\mu_c)$
    \end{enumerate}
    を満たすこと。
  \end{definition}

\end{frame}

\begin{frame}{モデル 3: 加点後スコアと score-feasibility}
  $c\in V_i$ に対する\textbf{加点後スコア}:
  \[
    \sigma_c(i):=
    \begin{cases}
      \sigma(i) & c=o,\\[2pt]
      \sigma(i)+k_t & c=t\in\T(i).
    \end{cases}
  \]
  留保カテゴリは共通の順序を持ち，資格・容量・選好だけが異なる（$k=0$ も含む）。

  % \begin{definition}[Cutoff と score-feasibility]
    \[
      e_o(\mu):=
      \begin{cases}
        \min_{i\in\mu_o}\sigma(i) & \text{if } \mu_o\neq\emptyset,\\
        -\infty & \text{if } \mu_o=\emptyset,
      \end{cases}
      \qquad
      e_t^+(\mu):=
      \begin{cases}
        \min_{i\in\mu_t}\sigma_t(i) & \text{if } \mu_t\neq\emptyset,\\
        +\infty & \text{if } \mu_t=\emptyset.
      \end{cases}
    \]
    feasible な $\mu$ が \textbf{score-feasible} とは
    $e_o(\mu)-e_t^+(\mu)\le\delta\ (\forall t\in T)$。
  % \end{definition}

  % \vspace{0.3em}
  % {\small $k$ と $\delta$ は別の役割: 加点は制度的に正当な補正，
  % $\delta$ は\emph{加点後に残る}差だけを縛る。
  % $k=0$ なら $\delta$ が生スコア差を縛り，これが 10 点ルールに一致する。}
\end{frame}

\begin{frame}{性質 1: Non-wastefulness}
  空席があっても，そこに入れると点数制約が破れるなら浪費とは呼ばない。

  \begin{definition}[Non-wastefulness]
    feasible な $\mu$ が \textbf{non-wasteful} とは，
    $i\in I\setminus M(\mu)$ と $c\in V_i$ で，
    $i$ を $c$ に加えた結果が
    \textbf{feasible かつ score-feasible} になるものが存在しないこと。
  \end{definition}

\end{frame}

\begin{frame}{性質 2: Within-Category Fairness}
  \begin{definition}[Within-category fairness]
    feasible な $\mu$ が \textbf{fair within categories} とは，
    次を全て満たす $t\in T$ と $i,j$ が存在しないこと:
    \[
      \text{(i) } t\in\T(i)\cap\T(j),\quad
      \text{(ii) } \sigma_t(i)>\sigma_t(j),\quad
      \text{(iii) } j\in\mu_t,\quad
      \text{(iv) } i\notin M(\mu).
    \]
  \end{definition}

  \vspace{0.3em}
  留保カテゴリのみを対象とするのは意図的で，
  これは 2 つのメカニズムが\textbf{共に}満たす最小の概念である。

  % \begin{block}{補題: 点数制約に相対化する必要はない}
  %   $j\in\mu_t$ を $\sigma_t(i)>\sigma_t(j)$ なる未割当 $i$ で\textbf{置き換える}と，
  %   $\mu'$ は feasible で，$\mu$ が score-feasible なら $\mu'$ も score-feasible。
  %   （$e_t^+$ は上がり，他は不変）
  % \end{block}

  % \vspace{0.2em}
  % {\small 非対称性: 新規に\emph{追加}すると cutoff は下がりうるが，
  % 上位者への\emph{置換}では下がらない。だから性質 1 だけが score 条項を持つ。}
\end{frame}

% \begin{frame}{性質 3: No-Justified-Envy}
%   カテゴリ版の公平性。\textbf{未割当}学生の不満に注目する。
%   \textbf{一般カテゴリ}への不満も許す点が性質 2 と異なる。

%   \begin{definition}[Within-category fairness]
%     feasible な $\mu$ が \textbf{fair within categories} とは，
%     次を満たす $i$ と $c\in V_i$ が存在しないこと:
%     \begin{itemize}
%       \item[(i)] $i\notin M(\mu)$，または $i\in\mu_{c'}$ かつ $c\succ_i c'$
%       \item[(ii)] $|\mu_c|<q^c$，または $\exists j\in\mu_c$ で $\sigma_c(j)<\sigma_c(i)$
%     \end{itemize}
%   \end{definition}

%   \vspace{0.3em}
%   \begin{itemize}
%     \item within-category fairness は留保カテゴリ内の順位違反だけを見る
%     \item ここで score 条項を省くのは\textbf{本質的な選択}:
%     (ii) は\emph{空席}の要求も許し（cutoff が下がる），
%     (i) は $c'=o$ からの離脱を許す（open cutoff が\emph{上がる}）ので，
%     補題は拡張されない
%   \end{itemize}

%   \vspace{0.3em}
%   {\small 構造的循環: feasible な留保割当は open cutoff に依存し，
%   open cutoff は留保割当後に誰が残るかに依存する。}
% \end{frame}

\begin{frame}{SPSE: Student-Proposing Score-Elevated Mechanism}
  加点後スコアの下での deferred acceptance。
  \begin{enumerate}
    \item $\mu_c\gets\emptyset$，全員未割当。$V_i=\{o\}\cup\T(i)$ と $\sigma_c(i)$ を計算
    \item 未応募カテゴリを残す未割当学生がいる間:
    \begin{itemize}
      \item $i$ は最も好む未応募の $c\in V_i$ に応募
      \item $c$ は $\mu_c\cup\{i\}$ のうち $\sigma_c$ 上位を高々 $q^c$ 人保留
      \item 拒否された学生は未割当に戻る
    \end{itemize}
  \end{enumerate}

  \vspace{0.4em}
  \begin{itemize}
    \item 学生は同時に高々一つのカテゴリにしか保留されない
    $\Rightarrow$ one-to-one counting を保つ
    \item 各学生は各カテゴリに高々 1 回応募するのでループは $|I|(|T|+1)$ 回以下で停止
    \item 通常の student-proposing DA との違いは，
    応募先が\emph{カテゴリ}であることと留保側の score elevation のみ
  \end{itemize}

  {\small 〔Gale and Shapley, 1962; Roth, 2008〕}
\end{frame}

\begin{frame}{SPSE の性質}
  \begin{center}
  \begin{tabular}{ll}
    \toprule
    性質 & 結果 \\
    \midrule
    Feasibility & 成立 \\
    Non-wastefulness & 成立 \\
    Within-category fairness & 成立 \\
    \midrule
    Score-feasibility & \textbf{保証されない} \\
    \bottomrule
  \end{tabular}
  \end{center}

  \vspace{0.5em}
  \begin{itemize}
    \item \textbf{Non-wastefulness}: 未割当 $i$ は $V_i$ の全カテゴリに応募済みで拒否されている。
    カテゴリは満員のときしか拒否せず，一度満員になれば容量を空けない $\Rightarrow$ $|\mu_c|=q^c$
    \item \textbf{Within-category fairness}: $j\in\mu_t$ より低い未割当 $i$ がいるなら，
    $i$ は $t$ に応募済みで拒否されているはずなので矛盾
  \end{itemize}
\end{frame}

\begin{frame}{SPSE は score-feasible とは限らない}
  \[
    T=\{t_1,t_2\},\quad q^o=q^{t_1}=q^{t_2}=1,\quad k=5,\quad \delta=10
  \]
  \begin{center}
  \begin{tabular}{c|c|c|c}
    \toprule
    学生 & 加点前スコア & 資格カテゴリ & 選好 \\
    \midrule
    $i_1$ & 100 & $\emptyset$ & $o$ \\
    $i_2$ & 95 & $\{t_1,t_2\}$ & $t_1\succ t_2\succ o$ \\
    $i_3$ & 82 & $\{t_2\}$ & $t_2\succ o$ \\
    \bottomrule
  \end{tabular}
  \end{center}

  \vspace{0.3em}
  第 1 ラウンドで各カテゴリが応募者を 1 人ずつ受け取り即停止:
  \[
    \mu_o=\{i_1\},\quad \mu_{t_1}=\{i_2\},\quad \mu_{t_2}=\{i_3\},\qquad
    e_o-e_{t_2}^+=100-87=13>10=\delta.
  \]

  一方 $\nu_o=\{i_1\}$, $\nu_{t_2}=\{i_2\}$, $\nu_{t_1}=\emptyset$ は
  score-feasible・non-wasteful・within-category fair（gap は $0$，空の $t_1$ は自明）。

  \vspace{0.3em}
  {\small $\nu$ が優れていると主張はしない: SPSE は 3 人，$\nu$ は 2 人を割り当て，
  $\nu$ 自身は別の意味で不満を残す。
  主張は「3 性質は score-feasibility を導かない」ことである。}
\end{frame}

% \begin{frame}{Score-infeasibility の原因}
%   \begin{block}{原因 1: overlap 不要（既知；Ayg\"un and Turhan 2023b の指摘）}
%     $T=\{t_1\}$, $q^o=q^{t_1}=1$, $k=5$, $\delta=10$,
%     $i_1(100,\emptyset)$, $i_2(80,\{t_1\})$。
%     SPSE は $e_o-e_{t_1}^+=100-85=15>\delta$。
%     \\ 唯一の有資格者が単に弱すぎるだけで，\textbf{割当の判断に誤りはない}。
%   \end{block}

%   \begin{block}{原因 2: overlap 固有（本研究）}
%     前ページの例では $i_2$ はどちらのカテゴリでも同じ加点後スコアを持ち，
%     \textbf{弱すぎる候補者は一人もいない}。
%     $i_2$ 一人ならどちらに入っても点数制約を満たす。
%     \\ SPSE は本人の選好で $t_1$ に送り，$t_2$ をはるかに弱い $i_3$ に残す。
%     \\ one-to-one counting の下では，候補者集合ではなく
%     \textbf{routing（どこに数えるか）}が制約を破っている。
%   \end{block}

%   \vspace{0.3em}
%   単一資格ではそもそも routing が存在しないので，原因 2 は生じない。
%   以降はこの原因 2 を扱う。
% \end{frame}



\begin{frame}{PA-ODR}
  open cutoff $e$ を固定し，二部グラフ $G(e)$ を
  \[
    \text{左}=\mu_o \text{ に入らない学生},\quad \text{右}=T,\qquad
    (i,t)\in G(e)\iff t\in\T(i)\ \wedge\ \sigma(i)+k\ge e-\delta
  \]
  で定める。共通加点の下では，学生は自分の全カテゴリで admissible か全く駄目かのどちらかだが，
  one-to-one counting が\textbf{どれを埋めるか}の選択を強いるのでグラフが必要になる。

  \begin{block}{\textsc{PrefLexMaxMatching}: どの辺を採用するか}
    \begin{enumerate}
      \item まず，できるだけ多くの留保席を埋める
      \item 同じ人数なら，より高スコアの学生を優先する
      \item それでも同じなら，学生のカテゴリ希望をできるだけ尊重する
      \item まだ同じ候補が複数あれば，その中から一つ選ぶ
    \end{enumerate}
  \end{block}

  \vspace{0.2em}
  \begin{itemize}
    \item \textbf{補題}: 基準 1--2 だけで\textbf{合格者集合}が一意に決まる
    （横断マトロイドの基＋相異なる重みに対する greedy）$\Rightarrow$ 基準 3 は\emph{どこに}置くかのみを決める
    \item 4 基準は 1 回の最大重み二部マッチングで実現でき，重みは $O(|I||T|)$ ビット $\Rightarrow$ 多項式時間
  \end{itemize}
\end{frame}

\begin{frame}{PA-ODR 2: アルゴリズム（閾値調整）}
  入力に $T$ 上の固定線形順序 $\prec$ を含める。
  \begin{enumerate}
    \item $\hat q^o\gets q^o$, $\hat q^t\gets q^t$
    \item \textbf{繰り返し}:
    \begin{enumerate}
      \item[(a)] 生スコア上位 $\min\{\hat q^o,|I|\}$ 人を $o$ に割り当て，$e\gets\min_{i\in\mu_o}\sigma(i)$
      \item[(b)] $G(e)$ を構成（$t$ の容量は $\hat q^t$）し，\textsc{PrefLexMaxMatching} で $(\mu_t)_{t\in T}$ を得る
      \item[(c)] $|\mu_t|<\hat q^t$ なる $t$ があれば，$\prec$ 最初のものを $t^\ast$ として
      \[
        \hat q^{t^\ast}\gets\hat q^{t^\ast}-1,\qquad \hat q^o\gets\hat q^o+1
      \]
      \item[(d)] なければ $\mu$ を出力
    \end{enumerate}
  \end{enumerate}

  \vspace{0.3em}
  \begin{itemize}
    \item 一般枠を先に生スコアで埋め，留保側は cutoff で feasible な組だけを対象にする
    \item 選好は\textbf{既に feasible な}留保配置の間の tie-break にしか使わない
    \item \textbf{停止規則は score-feasibility が要求する以上に強い}:
    埋まらない留保席があれば必ず de-reserve する（点数規則が破れる場合に限らない）
    \item \textbf{単調性}: $e^{(1)}\ge e^{(2)}\ge\cdots$，$\mu_o^{(1)}\subseteq\mu_o^{(2)}\subseteq\cdots$。
    $e$ が下がると辺は増えるが，同時に左側の頂点が 1 つ減る（どちらも支配しない）
  \end{itemize}
\end{frame}

\begin{frame}{PA-ODR 3: 動作例}
  \[
    T=\{t_1,t_2\},\quad q^o=q^{t_1}=q^{t_2}=1,\quad k=5,\ \delta=10,\qquad
    i_1(100,\emptyset),\ i_2(84,\{t_1\}),\ i_3(83,\{t_2\})
  \]
  1 回目: $\mu_o=\{i_1\}$, $e=100$。admissible には $\sigma(i)+5\ge 90$ が必要だが
  $89,88$ とも届かず，$G(100)$ は辺を持たない。

  \vspace{0.4em}
  \begin{columns}[t]
    \begin{column}{0.48\textwidth}
      \textbf{$t_1\prec t_2$ のとき}
      \begin{itemize}\small
        \item $t_1$ を de-reserve: $\hat q=(2,0,1)$
        \item $\mu_o=\{i_1,i_2\}$, $e=84$，条件は $\sigma(i)+5\ge 74$
        \item $i_3$ が届き $(i_3,t_2)$ が出現
        \item 出力 $\mu_o=\{i_1,i_2\}$, $\mu_{t_2}=\{i_3\}$
        \\ \textbf{de-reserve は 1 席}
      \end{itemize}
    \end{column}
    \begin{column}{0.48\textwidth}
      \textbf{$t_2\prec t_1$ のとき}
      \begin{itemize}\small
        \item $t_2$ を de-reserve: $\hat q=(2,1,0)$
        \item $\mu_o=\{i_1,i_2\}$, $e=84$
        \item $t_1$ の唯一の有資格者 $i_2$ が open に吸収され，$t_1$ への辺が無い
        \item 2 回目の de-reservation が起き $\mu_o=\{i_1,i_2,i_3\}$
        \\ \textbf{de-reserve は 2 席}
      \end{itemize}
    \end{column}
  \end{columns}

  % \vspace{0.4em}
  % どちらも 3 人を割り当てるので差はカバレッジではなく\textbf{制度的}:
  % 一方では留保枠が生き残り $i_3$ は $t_2$ 経由で入るが，
  % 他方では両枠が消え $i_3$ は $o$ 経由になる。
  % \\ {\small 注: これは overlap 由来ではなく cutoff のフィードバック由来（単一資格でも起きる）。}
\end{frame}

\begin{frame}{PA-ODR 4: 性質}
  \begin{block}{命題（停止性・多項式時間・feasibility・score-feasibility・non-wastefulness）}
    \begin{itemize}
      \item \textbf{停止}: 非終了反復ごとに $\sum_t\hat q^t$ が 1 減るので高々 $\sum_t q^t$ 回
      \item \textbf{多項式時間}: 先に $q^c\gets\min\{q^c,|I|\}$ と正規化すれば反復は $O(|I||T|)$ 回
      \item \textbf{Score-feasibility}: 留保割当は全て $G(e)$ の辺に沿うので
      $\sigma(i)+k\ge e-\delta$，よって $e_t^+(\mu)\ge e_o(\mu)-\delta$
      \item \textbf{Non-wastefulness（最終容量 $\hat q$ に関して）}:
      停止時は $|\mu_t|=\hat q^t$ が全 $t$ で成立し，$o$ も上位 $\hat q^o$ 人で満員
    \end{itemize}
  \end{block}

  \begin{block}{命題（within-category fairness)}
    未割当 $i$ が $j\in\mu_t$ より $\sigma_t$ で上位なら，$(i,t)$ は自動的に $G(e)$ の辺。
    $j$ を $i$ に置換するとマッチ数は同じでスコアベクトルが厳密に改善し，
    基準 1--2 の最適性に矛盾。
  \end{block}

  % \vspace{0.2em}
  % {\small SPSE と違い，性質の源泉は拒否の単調性ではなく
  % \emph{最終 feasibility graph の中でのマッチングの最適性}である。}
\end{frame}

% \begin{frame}{PA-ODR が諦めるもの}
%   いずれも同じ源: \textbf{一般カテゴリを選好を見ずに生スコアだけで先に埋める}こと。

%     $T=\{t_1\}$, $q^o=q^{t_1}=1$, $k=5$, $\delta=10$,
%     $i_1(100,\{t_1\})$, $i_2(90,\{t_1\})$，共に $t_1\succ o$。
%     \\ PA-ODR: $\mu_o=\{i_1\}$, $e=100$, $95\ge 90$ より $\mu_{t_1}=\{i_2\}$（de-reserve なし，score-feasible）。
%     \\ だが $i_1$ は $t_1$ を好み $\sigma_{t_1}(i_1)=105>95$。
%     請求を認めた $\mu_{t_1}=\{i_1\},\mu_o=\{i_2\}$ も score-feasible なので\textbf{点数規則は言い訳にならない}。
%     SPSE はまさにこちらを返す。
%   \end{block}

%     $i_1(100,\emptyset)$, $i_2(95,\{t_1,t_2\})$, $i_3(94,\{t_1,t_2\})$，
%     $t_1\succ_{i_2}t_2\succ_{i_2}o$, $o\succ_{i_3}t_1\succ_{i_3}t_2$。
%     \\ 正直申告では順位ベクトルが $\alpha:(3,1)$, $\beta:(2,2)$ で基準 3 は $\beta$ を選び $i_2$ は $t_2$。
%     $t_1\succ o\succ t_2$ と偽ると $\alpha:(3,1)$, $\beta:(3,2)$ となり $\alpha$ が選ばれ $i_2$ は $t_1$ を得る。
%   \end{block}

%   \vspace{0.2em}
%   さらに: de-reservation は\textbf{順序依存で最小でない}。
%   ただし最小 de-reservation 問題 \textsc{MinDeres} は
%   $O(\prod_t(q^t+1)\cdot\Phi)$ の列挙で解け，$|T|$ が定数なら多項式時間
%   $\Rightarrow$ 順序依存は目的関数の難しさではなく\emph{貪欲なスケジュール}の性質。
% \end{frame}

\begin{frame}{2 つのメカニズムの比較}
  \begin{center}
  \small
  \begin{tabular}{lll}
    \toprule
    \textbf{性質} & \textbf{SPSE} & \textbf{PA-ODR} \\
    \midrule
    One-to-one counting & 成立 & 成立 \\
    Feasibility & 成立 & 成立（最終 $\hat q$ に関して） \\
    Non-wastefulness & 成立 & 成立（最終 $\hat q$ に関して） \\
    Within-category fairness & 成立 & 成立 \\
    \midrule
    Score-feasibility & 保証されない & 構成的に保証 \\
    容量の再配分 & $q$ に固定 & 内生的・最小でない \\
    % カテゴリ選好の利用 & 応募を通じて & 辞書式 tie-break を通じて \\
    \bottomrule
  \end{tabular}
  \end{center}

  % \vspace{0.4em}
  % 横線より上は両者とも満たす（PA-ODR は自ら生成した実効容量に関して）。
  % 下では真っ向から異なる。
\end{frame}

\begin{frame}{まとめ}
  \begin{itemize}
    \item SPSE は within-category fairness を満たすが score-feasible でない。
    \item PA-ODR は点数制約を hard に扱い，within-category fairness を弱めずに保つが，
    空席のままで済んだ枠まで retire する
    % \item 固定容量と score-feasibility の緊張は特定アルゴリズムの人工物ではない
  \end{itemize}

\end{frame}

% \begin{frame}{相談事項 1: 現状の新規性と懸念}
%   \begin{block}{overlap に固有と言い切れる主張は 2 つだけ}
%     \begin{enumerate}
%       \item \textbf{Routing}: 複数資格の学生をどのカテゴリで数えるかという判断だけで
%       score-feasibility が決まりうる（候補者の弱さとは独立の原因）
%       \item \textbf{Hall 型ギャップ}: 各枠が個別には score-feasible に充足可能でも同時には不可能。
%       $N(\{t\})$ が互いに素なら消えるので，差はちょうど overlap
%     \end{enumerate}
%   \end{block}

%   \vspace{0.3em}
%   一方で，\textbf{overlap がなくても出る}部分がある:
%   \begin{itemize}
%     \item score-elevated 政策が cutoff 規則を破ること自体
%     〔Ayg\"un and Turhan, 2023b〕
%     \item de-reservation の順序依存・非最小性（cutoff のフィードバック由来）
%   \end{itemize}

%   \vspace{0.4em}
%   懸念: 現状の貢献は\textbf{既知の枠組みへの観察と反例}に寄っており，
%   「確実に新しい」と言える定理が \textbf{Proposition（不可能性）1 本}に集中している。
%   \\ $\Rightarrow$ もう一段、\textbf{確実な新規結果}を足したい。
% \end{frame}

\begin{frame}{これからの方針}

\begin{itemize}
    \item カテゴリを跨いだ公平性を何らかの形で考えたい
    \item DAベースではない何らかのアルゴリズムを考えたい
\end{itemize}

\end{frame}

\begin{frame}{References}
  \footnotesize
  \begin{itemize}
    \item Gale and Shapley (1962), College admissions and the stability of marriage.
    \item Roth (2008), Deferred acceptance algorithms: History, theory, practice, and open questions.
    \item Ehlers, Hafalir, Yenmez, and Yildirim (2014), School choice with controlled choice constraints.
    \item Kominers and S\"onmez (2016), Matching with slot-specific priorities.
    \item Goto, Kojima, Kurata, Tamura, and Yokoo (2017), Designing matching mechanisms under general distributional constraints.
    \item Kurata, Hamada, Iwasaki, and Yokoo (2017), Controlled school choice with soft bounds and overlapping types.
    \item S\"onmez and Yenmez (2020; 2022), Affirmative action with overlapping reserves; Affirmative action in India via vertical, horizontal, and overlapping reservations.
    \item Ayg\"un and Turhan (2020), Dynamic reserves in matching markets.
    \item Ayg\"un and B\'o (2021), College admission with multidimensional privileges: The Brazilian affirmative action case.
    \item Ayg\"un and Turhan (2023a), How to de-reserve reserves: Admissions to technical colleges in India.
    \item Ayg\"un and Turhan (2023b), Designing direct matching mechanisms for India with non-transferable private colleges（10 点ルール）.
  \end{itemize}
\end{frame}

% ============================================================================
% 以下は旧版スライドのアーカイブ（コメントアウトのまま保持）
% ============================================================================

% --- 旧「SPSE の性質」表の行 ---------------------------------------------
%    Across-category fairness & \textbf{不成立（反例あり）} \\

% --- 旧「PA-ODR で評価する性質」の項目 -------------------------------------
%    \item PA-ODR は人数最大 $\to$ スコアベクトル辞書式最大の基準で，この fairness notion と整合的に選ぶ
%    \item score-feasibility は hard constraint として別に保証する

% --- 旧「PA-ODR の結果」表の行 --------------------------------------------
%    制約付き across-category fairness & 成立 \\
%    強い stability & \textbf{一般には不成立} \\
%    容量再配分 & 成立 \\

% \begin{frame}{性質 4: Across-Category Fairness}
%   資格を共有しない $i,j$ の間で，$i$ を入れる再配置が
%   \textbf{割り当てられている学生のスコアベクトルを改善する}のに $i$ が落ちている状況を禁止する。

%   \begin{block}{悪い状況（unfair）：次を\textbf{すべて}満たす対 $(i,j)$ が存在}
%     \begin{itemize}
%       \item[(1)] $c=\mu(j)\notin V_i$（$i$ は $j$ のカテゴリに割り当てられない。必ず $t\notin\T(i)$）
%       \item[(2)] $j$ を外し $i$ を収容する score-feasible な再配置 $\mu'$
%                  （他の全員は維持，$\bigl(M(\mu)\setminus\{j\}\bigr)\cup\{i\}\subseteq M(\mu')$）で，
%                  \textbf{加点済みのスコアベクトル が $\mu$ より辞書式に大きい}ものが存在
%     \end{itemize}
%   \end{block}

%   \begin{definition}[Across-category fairness]
%     $\mu$ が \textbf{fair across categories} であるとは，
%     $(1)$--$(2)$ を満たす対 $(i,j)$ がひとつも存在しないことをいう。
%   \end{definition}

%   {\footnotesize 比較は生スコアでなく加点後スコアで行う。
%   PA-ODR の選択基準（人数最大 $\to$ elevated 辞書式最大）と整合する。}
% \end{frame}

% \begin{frame}{スコアベクトルとは}
%   あるマッチング $\mu$ に対して，
%   割り当てられている全学生の \textbf{加点後スコア} を大きい順に並べた列を
%   スコアベクトルと呼ぶ。

%   \[
%     s(\mu)=\bigl(s_1(\mu),s_2(\mu),\dots,s_{|M(\mu)|}(\mu)\bigr),
%     \qquad
%     s_1(\mu)\ge s_2(\mu)\ge \cdots
%   \]

%   ここで各学生 $i\in M(\mu)$ の値は，
%   実際に割り当てられたカテゴリ $\mu(i)$ における 加点後
%   \[
%     \sigma_{\mu(i)}(i)
%   \]
%   である。

%   \vspace{0.5em}
%   例：
%   \[
%     \mu=(A\to o,\ B\to t_1,\ C\to t_2),\quad
%     \bigl(\sigma_o(A),\sigma_{t_1}(B),\sigma_{t_2}(C)\bigr)=(100,95,90)
%   \]
%   なら
%   \[
%     s(\mu)=(100,95,90).
%   \]

%   PA-ODR では同じ人数の候補の中で，このベクトルを\textbf{辞書式}に最大化する。
% \end{frame}


% \begin{frame}{性質 5: Score-feasibility}
%   マッチング $\mu$ が score-feasible であるとは，
%   全ての $t\in T$ について，
%   $\mu_o\neq\emptyset$ かつ $\mu_t\neq\emptyset$ なら
%   \[
%     e_o(\mu)-e_t^+(\mu)\le \delta
%   \]
%   を満たすことをいう。

%   \vspace{0.5em}
%   ここで
%   \[
%     e_o(\mu)=\min_{i\in\mu_o}\sigma(i),
%     \qquad
%     e_t^+(\mu)=\min_{i\in\mu_t}\sigma_t(i)
%   \]
%   である。

%   {\small [Ayg\"un and Turhan, 2020; Ayg\"un and Turhan, 2023]}
% \end{frame}

% \begin{frame}{PA-ODR の証明 3}
%   \textbf{制約付き across-category fairness}
%   \begin{itemize}
%     \item 反対に，PA-ODR の結果 $\mu$ が unfair だと仮定する
%     \item すると定義より，ある $(i,j)$ について，
%     $j$ を外して $i$ を入れる score-feasible な再配置 $\mu'$ が存在し，
%     \[
%       s(\mu')>_{\mathrm{lex}} s(\mu)
%     \]
%     が成り立つ
%     \item しかも $\mu'$ は $|M(\mu')|=|M(\mu)|$ を満たすので，
%     PA-ODR が比較している「同じ人数の score-feasible 候補」のひとつである
%     \item しかし PA-ODR はその集合の中で辞書式最大のスコアベクトルを選ぶので矛盾
%   \end{itemize}

%   \vspace{0.4em}
%   よって PA-ODR は，この制約付きの意味では across-category fairness を満たす。

%   \vspace{0.4em}
%   ただし，割り当て済み学生がより好むカテゴリの在籍者を羨む
%   \textbf{stability} まで要求すると，一般には成立しない。
% \end{frame}



% \begin{frame}{反例：SPSE も PA-ODR も across-category fairness を破る}
%   \small
%   設定：$A$(open, $\sigma=200$)，$k\in\{t_1,t_2\}$($\sigma=100$, 選好 $t_1\succ t_2$)，
%   $i\in\{t_1\}$($\sigma=90$)，$j\in\{t_2\}$($\sigma=80$)，
%   $k_{t_1}=30,\ k_{t_2}=0$，$\delta$ 大。$A\to o$ で $e$ は十分高い。

%   \begin{center}
%   \begin{tabular}{lcc}
%     \toprule
%     留保割当 & elevated vector & 人数 \\
%     \midrule
%     $\mu=(k\to t_1,\ j\to t_2)$ & $\{130,\,80\}$ & 2 \\
%     $\mu'=(i\to t_1,\ k\to t_2)$ & $\{120,\,100\}$ & 2 \\
%     \bottomrule
%   \end{tabular}
%   \end{center}

%   \begin{itemize}
%     \item SPSE：$t_1$ で $k$(130)$>$$i$(120) が保留され，$i$ は $o$ で $A$ に負け未マッチ
%     \item PA-ODR：elevated vector 辞書式最大で $130>120$ より $\mu$ を採用
%     \item どちらも出力は $\mu$。だが $\mu'$ は score-feasible で人数も同じ
%           $\Rightarrow$ raw 上位の $i$（90）が $j$（80）を跨いで落ちており unfair
%   \end{itemize}

%   \begin{block}{結論}
%     raw-merit の across-category fairness は，elevated score で優先する
%     score-elevated メカニズム（SPSE・PA-ODR）では一般に達成できない。
%   \end{block}
% \end{frame}

% \begin{frame}{動作例 1}
%   設定:
%   \[
%     q^o=q^{t_1}=q^{t_2}=1,\quad
%     \delta=10,\quad
%     k_{t_1}=5,\ k_{t_2}=5
%   \]

%   \begin{center}
%   \begin{tabular}{c|c|c|c}
%     \toprule
%     学生 & 加点前スコア & 資格カテゴリ & カテゴリ選好 \\
%     \midrule
%     $A$ & 100 & $\emptyset$ & $o$ \\
%     $B$ & 85 & $\{t_1,t_2\}$ & $t_2\succ_B t_1\succ_B o$ \\
%     $C$ & 85 & $\{t_1\}$ & $t_1\succ_C o$ \\
%     $D$ & 85 & $\{t_2\}$ & $t_2\succ_D o$ \\
%     \bottomrule
%   \end{tabular}
%   \end{center}
% \end{frame}

% \begin{frame}{動作例 1: 辺の構成}
%   open は 加点前スコア 最大の $A$ を選ぶので
%   \[
%     \mu_o=\{A\},\qquad e=100.
%   \]

%   留保辺は
%   \[
%     \sigma(i)+k_t\ge 100-10=90
%   \]
%   を満たすときに張られる。

%   \begin{itemize}
%     \item $B$: $85+5=90$ なので $t_1,t_2$ の両方に届く
%     \item $C$: $85+5=90$ なので $t_1$ に届く
%     \item $D$: $85+5=90$ なので $t_2$ に届く
%   \end{itemize}

%   人数最大・スコアも最大の割当は
%   \[
%     (B,t_1),(D,t_2)
%     \qquad\text{または}\qquad
%     (C,t_1),(B,t_2)
%   \]
%   の 2 つある（スコアはどちらも同じ）。

%   それぞれで，割り当てられた学生が\textbf{何番目の希望}に入るかを見る：
%   \begin{itemize}
%     \item $(B,t_1),(D,t_2)$：$B$ は第2希望，$D$ は第1希望
%     \item $(C,t_1),(B,t_2)$：$C$ も $B$ も第1希望
%   \end{itemize}
%   後者の方が希望に近いので，PA-ODR は $(C,t_1),(B,t_2)$ を選ぶ。
% \end{frame}

% \begin{frame}{動作例 1: 性質確認}
%   最終割当:
%   \[
%     \mu_o=\{A\},\qquad \mu_{t_1}=\{C\},\qquad \mu_{t_2}=\{B\}.
%   \]

%   cutoff は
%   \[
%     e_o=100,\qquad e_{t_1}^+=90,\qquad e_{t_2}^+=90.
%   \]
%   したがって
%   \[
%     100-90=10\le 10,\qquad 100-90=10\le 10.
%   \]

%   \begin{itemize}
%     \item score-feasibility はちょうど境界で成立する
%     \item $B$ は $t_1,t_2$ のどちらでも feasible だが，選好により $t_2$ に入る
%     \item つまりこの例ではカテゴリ選好が実際に tie-break として働く
%   \end{itemize}
%   よって score-feasibility を満たす範囲でカテゴリ選好が反映される。
% \end{frame}

% \begin{frame}{新しい論点: 強い stability との衝突}
%   PA-ODR は人数最大化を優先するため，
%   割り当て済み学生のカテゴリ選好まで含めた stability とトレードオフの関係にある

%   \vspace{0.5em}
%   \begin{block}{強い stability の考え方}
%     学生 $i$ が現在のカテゴリより $c$ を好み，
%     $c$ の在籍者 $j$ より $c$ で高い 加点後スコア を持つなら，
%     $(i,j)$ は blocking pair になりうる。
%   \end{block}

%   \vspace{0.5em}
%   これは，未割当学生だけを見る fairness より強い要求である。
%   この強い要求を入れると，
%   \[
%     \text{人数最大化}
%     \quad+\quad
%     \text{score-feasibility}
%     \quad+\quad
%     \text{stability}
%   \]
%   を同時に保証することは一般に難しい。
% \end{frame}

% \begin{frame}{非両立の小例}
%   \small
%   \[
%     q^o=q^{t_1}=q^{t_2}=1,\quad
%     k_{t_1}=5,\quad k_{t_2}=10,\quad \delta=10
%   \]

%   \begin{center}
%   \begin{tabular}{c|c|c|c}
%     \toprule
%     学生 & 加点前スコア & 資格カテゴリ & 選好 \\
%     \midrule
%     $A$ & 100 & $\emptyset$ & $o$ \\
%     $B$ & 86 & $\{t_1,t_2\}$ & $t_2\succ_B t_1\succ_B o$ \\
%     $C$ & 84 & $\{t_1\}$ & $t_1\succ_C o$ \\
%     $D$ & 82 & $\{t_2\}$ & $t_2\succ_D o$ \\
%     \bottomrule
%   \end{tabular}
%   \end{center}

%   関係する 加点後スコア は
%   \[
%     \sigma_{t_1}(B)=91,\quad
%     \sigma_{t_2}(B)=96,\quad
%     \sigma_{t_1}(C)=89,\quad
%     \sigma_{t_2}(D)=92.
%   \]

%   score-feasible で人数最大の割当は
%   \[
%     \mu^\ast=(A\to o,\ B\to t_1,\ D\to t_2)
%   \]
%   で，人数は $3$。
% \end{frame}

% \begin{frame}{非両立の直観}
%   $\mu^\ast=(A\to o,\ B\to t_1,\ D\to t_2)$ では
%   \[
%     e_o=100,\quad e_{t_1}^+=91,\quad e_{t_2}^+=92
%   \]
%   なので score-feasible である。

%   \vspace{0.5em}
%   しかし $B$ は
%   \[
%     t_2\succ_B t_1
%     \qquad\text{かつ}\qquad
%     \sigma_{t_2}(B)=96>92=\sigma_{t_2}(D)
%   \]
%   を満たす。

%   \vspace{0.5em}
%   つまり，人数を $3$ に保つためには $B$ を $t_1$ に置く必要があるが，
%   その結果 $B$ は $t_2$ の $D$ を羨む。

%   \begin{block}{含意}
%     PA-ODR は人数最大化と score-feasibility を取る代わりに，
%     強い stability を犠牲にする場合がある。
%   \end{block}
% \end{frame}

% --- 旧「動作例（A,B,C,D による de-reservation の 2 回反復）」（原文のまま保持）---
% \begin{frame}{動作例}
%   設定:
%   \[
%     q^o=q^{t_1}=q^{t_2}=1,\quad
%     \delta=10,\quad
%     k=5
%   \]
% 
%   \begin{center}
%   \begin{tabular}{c|c|c|c}
%     \toprule
%     学生 & 加点前スコア & 資格カテゴリ & カテゴリ選好 \\
%     \midrule
%     $A$ & 100 & $\emptyset$ & $o$ \\
%     $B$ & 92 & $\{t_1\}$ & $t_1\succ_B o$ \\
%     $C$ & 80 & $\{t_2\}$ & $t_2\succ_C o$ \\
%     $D$ & 78 & $\{t_1,t_2\}$ & $t_2\succ_D t_1\succ_D o$ \\
%     \bottomrule
%   \end{tabular}
%   \end{center}
% \end{frame}
% 
% \begin{frame}{動作例: 1回目}
%   初回は
%   \[
%     \mu_o=\{A\},\qquad e=100.
%   \]
%   留保辺条件は
%   \[
%     \sigma(i)+k\ge 90.
%   \]
% 
%   \begin{itemize}
%     \item $B$: $92+5=97$ なので $t_1$ に届く
%     \item $C$: $80+5=85<90$ なので $t_2$ に届かない
%     \item $D$: $78+5=83<90$ なのでどちらにも届かない
%   \end{itemize}
% 
%   したがって $t_2$ は埋まらず，
%   \[
%     \hat q^{t_2}\gets 0,\qquad \hat q^o\gets 2
%   \]
%   と de-reserve される。
% 
%   \vspace{0.3em}
%   この時点で $D$ は $t_2\succ_D t_1$ だが，
%   そもそも feasible な辺を持たない。
% \end{frame}
% 
% \begin{frame}{動作例: 2回目}
%   2回目は
%   \[
%     \hat q^o=2,\qquad \hat q^{t_1}=1,\qquad \hat q^{t_2}=0.
%   \]
%   open は $A,B$ を選ぶので
%   \[
%     \mu_o=\{A,B\},\qquad e=92.
%   \]
% 
% 
%   今度は
%   \[
%     \sigma(D)+k=78+5=83\ge 82
%   \]
%   なので $D$ が $t_1$ に届く。
% 
%   一方で $t_2$ はすでに de-reserve されているので，
%   $D$ の第一希望 $t_2$ は選べない。
% 
%   最終的に
%   \[
%     \mu_o=\{A,B\},\qquad \mu_{t_1}=\{D\},\qquad \mu_{t_2}=\emptyset.
%   \]
% 
%   \begin{itemize}
%     \item $t_2$ は score 制約の下で埋められないので open に戻る
%     \item つまりカテゴリ選好は，de-reservation と score-feasibility を満たす範囲でだけ反映される
%   \end{itemize}
% \end{frame}
% 
% \begin{frame}{性質}
%   最終容量は
%   \[
%     \hat q^o=2,\qquad \hat q^{t_1}=1,\qquad \hat q^{t_2}=0
%   \]
%   であり，最終割当は
%   \[
%     \mu_o=\{A,B\},\qquad \mu_{t_1}=\{D\},\qquad \mu_{t_2}=\emptyset.
%   \]
% 
%   \begin{itemize}
%     \item Score-feasibility:
%     \[
%       e_o=92,\qquad e_{t_1}^+=83,\qquad 92-83=9\le 10
%     \]
%     \item Non-wastefulness w.r.t.\ $\hat q$:
%     $\hat q^{t_2}=0$ なので $t_2$ の空きカテゴリはもう存在しない。$o,t_1$ も満員なので，未割り当て学生を追加できない
%   \end{itemize}
%   
% \end{frame}

\end{document}
